% ch11 译文：原文第 6 节 6b（main.tex 第 3658–4215 行）
% [原文 p.43]
\begin{dfn}\label{dfn:GEquivQCohdg}
与 $\mathfrak{U}$ 相关的 \emph{$X$ 的 $G$-等变凝聚导出 dg-范畴}是 dg-范畴
$$
\Coh_G(X;\mathfrak{U})_{\dg},
$$
其定义方式与定义~\ref{dfn:dg-enhancement-Coh} 类似，只是把 $\check{C}^{\bullet}(\mathfrak{U},\mathcal{F}^{\bullet})$ 换成 $G$-不变的 \v{C}ech 上链复形 $\check{C}^{\bullet}_G(\mathfrak{U},(\mathcal{F}^{\bullet},\theta))$。
\end{dfn}

直接验证表明，$\Coh_G(X;\mathfrak{U})_{dg}$ 是 dg 范畴，并且是拟凝聚层的导出范畴 $D^b(\Coh_G(X))$ 的 dg 加强。

\begin{yizhu}[dg 加强到底加强了什么]
把 $\Coh_G(X;\mathfrak{U})_{\dg}$ 称作 $D^b(\Coh_G(X))$ 的「dg 加强」，意思是：前者是一个 dg 范畴，其第零阶同伦范畴 $H^0$ 与三角范畴 $D^b(\Coh_G(X))$ 等价，同时每个 $\hom$ 都保留为一条上链复形，因而记下了导出 $\Hom$ 的全部链层面信息。只有在这层加强之上，才谈得上 dg-函子、拟等价以及 $\widehat{G}$ 的严格作用；若只剩下三角范畴，这些结构都会丢失。原文此处「quasi-coherent sheaves」与 $D^b(\Coh_G(X))$ 的字面含义略有出入，译文照原文保留。
\end{yizhu}

\begin{rmk}
要求 $X$ 存在由 $G$-不变仿射开集组成的开覆盖，是一个略带技术性的假设。由 Sumihiro \cite{Sumihiro74} 的工作，当 $X$ 正规、$G$ 为连通线性代数群时这样的覆盖存在，这已足以涵盖本文讨论的情形。
\end{rmk}

注意，$\Coh_G(X)$ 上的严格 $\widehat{G}$-作用
$$
\phi:\widehat{G}\rightarrow\Aut_{\str}(\Coh_G(X))
$$
自然诱导 $\Coh_G(X;\mathfrak{U})_{\dg}$ 上的严格 $\widehat{G}$-作用
$$
\phi:\widehat{G}\rightarrow\Aut_{\str}(\Coh_G(X;\mathfrak{U})_{\dg})
$$
（注~\ref{rmk:EmbeddingOfCochainComplexes}）。

对闭子群概形 $H\subseteq G$，限制函子 $\Res^G_H:\Coh_G(X)\rightarrow\Coh_H(X)$（定义~\ref{defn:RestrictionOfEquivariantSheaves}）诱导 dg-函子 $\Res_G^H:\Coh_G(X;\mathfrak{U})_{dg}\rightarrow\Coh_H(X;\mathfrak{U})_{dg}$，并且我们有受限的严格 $H^\perp$-作用
$$
\phi:H^\perp\rightarrow\Aut_{\str}(\Coh_G(X)_{\dg}).
$$

\begin{yizhu}[限制函子的记号方向]
原文对同一个限制 dg-函子用了两种写法：正文写成 $\Res^G_H$，而推论~\ref{crl:RestrictionFunctionEm} 及其证明写成 $\Res_H^G$，上下标正好互换。这是原文记号不统一，译文一律照录、不作改动。读者只需记住：无论哪种写法，它都是 $\Coh_G(X;\mathfrak{U})_{\dg}\rightarrow\Coh_H(X;\mathfrak{U})_{\dg}$，即把 $G$-等变结构遗忘为 $H$-等变结构。
\end{yizhu}

% [原文 p.44]
\begin{defn}[限制 dg-函子]
设 $H\subseteq G$ 为闭子群概形，回顾限制函子 $\Res^G_H:\Coh_G(X)\rightarrow\Coh_H(X)$（定义~\ref{defn:RestrictionOfEquivariantSheaves}）。设存在 $X$ 的仿射开覆盖 $\mathfrak U:=(U_i)_{i\in I}$，使得每个 $U_i$ 都是 $G$-不变的。于是每个 $U_i$ 自动是 $H$-不变的。我们定义 dg-范畴之间的限制 dg-函子
$$
\Res_G^H:\Coh_G(X;\mathfrak{U})_{\dg}\rightarrow\Coh_H(X;\mathfrak{U})_{\dg},
$$
它在对象与态射上的映射如下：
\begin{itemize}
\item
对象：
$
(\mathcal{F}^{\bullet},\theta^{\bullet})
\mapsto
(\mathcal{F}^{\bullet},\theta^{\bullet}|_H)
$
其中
$$
(\mathcal{F}^{\bullet},\theta^{\bullet}|_H)
:=
\left(
\bigoplus_{\objectdegree\in\mathbb Z}
(\mathcal F^{\objectdegree},\theta^{\objectdegree}|_H),
d_{\mathcal F^\bullet}
:=
\left(d_{\mathcal{F}^{\bullet}}^i:\mathcal{F}^i\rightarrow\mathcal{F}^{i+1}\right)_{i\in\mathbb{Z}}
\right).
$$
这里每个 $(\mathcal{F}^i,\theta^i|_H)$ 是 $(\mathcal{F}^i,\theta^i)$ 在限制函子 $\Res^G_H:\Coh_G(X)\rightarrow\Coh_H(X)$ 下的像，而 $d_{\mathcal{F}^{\bullet}}^i\in\Hom_{\Coh_H(X)}(\mathcal F^i,\mathcal F^{i+1})$ 是态射 $d_{\mathcal{F}^{\bullet}}^i\in\Hom_{\Coh_G(X)}(\mathcal F^i,\mathcal F^{i+1})$ 在该函子下的像。
\item
态射：回顾 hom 空间的定义
\[
\resizebox{\linewidth}{!}{$\displaystyle
\hom_{\Coh_G(X;\mathfrak U)_{\dg}}^{\homdegree}
\left(
(\mathcal F^\bullet,\theta^\bullet),
(\mathcal F'^\bullet,\theta'^\bullet)
\right)
=
\bigoplus_{p+q=\homdegree}
\prod_{i_0<\cdots<i_p}
\bigoplus_{\objectdegree\in\mathbb Z}
\Hom_{\Coh_G(U_{i_0\cdots i_p})}
\left(
(\mathcal F^{\objectdegree},\theta^{\objectdegree})
|_{U_{i_0\cdots i_p}},
(\mathcal F'^{\objectdegree+q},\theta'^{\objectdegree+q})
|_{U_{i_0\cdots i_p}}
\right).$}
\]
每个直和项都有自然的限制
\[
\resizebox{\linewidth}{!}{$\displaystyle
\Hom_{\Coh_G(U_{i_0\cdots i_p})}
\left(
(\mathcal F^{\objectdegree},\theta^{\objectdegree})
|_{U_{i_0\cdots i_p}},
(\mathcal F'^{\objectdegree+q},\theta'^{\objectdegree+q})
|_{U_{i_0\cdots i_p}}
\right)
\rightarrow
\Hom_{\Coh_H(U_{i_0\cdots i_p})}
\left(
(\mathcal F^{\objectdegree},\theta^{\objectdegree}|_H)
|_{U_{i_0\cdots i_p}},
(\mathcal F'^{\objectdegree+q},\theta'^{\objectdegree+q}|_H)
|_{U_{i_0\cdots i_p}}
\right),$}
\]
它由函子 $\Res^G_H:\Coh_G(U_{i_0\cdots i_p})\rightarrow\Coh_H(U_{i_0\cdots i_p})$ 给出。因此有自然的限制
$$
\hom_{\Coh_G(X;\mathfrak U)_{\dg}}^{\homdegree}
\left(
(\mathcal F^\bullet,\theta^\bullet),
(\mathcal F'^\bullet,\theta'^\bullet)
\right)
\rightarrow
\hom_{\Coh_H(X;\mathfrak U)_{\dg}}^{\homdegree}
\left(
(\mathcal F^\bullet,\theta^\bullet),
(\mathcal F'^\bullet,\theta'^\bullet)
\right).
$$
\end{itemize}
\end{defn}

\begin{crl}
\label{crl:RestrictionFunctionEm}
设 $G$ 是可对角化群概形，作用在概形 $X$ 上，$H\subseteq G$ 是闭子群概形。设 $\mathfrak U:=(U_i)_{i\in I}$ 是 $X$ 的仿射开覆盖，使得每个 $U_i$ 都是 $G$-不变的（从而是 $H$-不变的）。\par
则限制 dg-函子
$
\Res_H^G:\Coh_G(X;\mathfrak{U})_{\dg}
\rightarrow
\Coh_H(X;\mathfrak{U})_{\dg}
$
诱导轨道范畴的拟嵌入
$$
\Coh_G(X;\mathfrak{U})_{\dg}/H^\perp
\hookrightarrow
\Coh_H(X;\mathfrak{U})_{\dg}.
$$
\end{crl}

\begin{proof}
设 $(\mathcal{F},\theta)$ 与 $(\mathcal{F}',\theta')$ 是 $\Coh_G(X;\mathfrak{U})_{\dg}$ 中的两个对象。
$$
\hspace{-80mm}
\hom_{\Coh_H(X)_{\dg}}
\left((\mathcal{F}^{\bullet},\theta^{\bullet}|_H),(\mathcal{F}'^{\bullet},\theta'^{\bullet}|_H)\right)
$$
$$
=
\displaystyle
\bigoplus_{p+q=\homdegree}
\prod_{i_0<\cdots<i_p}
\bigoplus_{\objectdegree\in\mathbb Z}
\Hom_{\Coh_H(U_{i_0\cdots i_p})}
\left(
(\mathcal F^{\objectdegree},\theta^{\objectdegree}|_H)
|_{U_{i_0\cdots i_p}},
(\mathcal F'^{\objectdegree+q},\theta'^{\objectdegree+q}|_H)
|_{U_{i_0\cdots i_p}}
\right)
$$
% [原文 p.45]
由命题~\ref{prop:HInvariantsOfHomSpace}，每个直和项分解为
$$
\hspace{-30mm}
\Hom_{\Coh_H(U_{i_0\cdots i_p})}
\left(
(\mathcal F^{\objectdegree},\theta^{\objectdegree}|_H)
|_{U_{i_0\cdots i_p}},
(\mathcal F'^{\objectdegree+q},\theta'^{\objectdegree+q}|_H)
|_{U_{i_0\cdots i_p}}
\right)
$$
$$
=
\bigoplus_{\lambda\in H^{\perp}}
\Hom_{\Coh_G(U_{i_0\cdots i_p})}
\left(
(\mathcal F^{\objectdegree},\theta^{\objectdegree})
|_{U_{i_0\cdots i_p}}(\lambda),
(\mathcal F'^{\objectdegree+q},\theta'^{\objectdegree+q}|_H)
|_{U_{i_0\cdots i_p}}
\right)
$$
另一方面，由 dg-范畴的轨道范畴的定义，
$$
\hspace{-80mm}
\hom_{\Coh_G(X)_{\dg}/H^{\perp}}
\left((\mathcal{F}^{\bullet},\theta^{\bullet}|_H),(\mathcal{F}'^{\bullet},\theta'^{\bullet}|_H)\right)
$$
$$
\hspace{-58mm}
=
\displaystyle
\bigoplus_{\lambda\in H^{\perp}}
\hom_{\Coh_G(X)_{\dg}}
\left((\mathcal{F}^{\bullet},\theta^{\bullet}|_H)(\lambda),(\mathcal{F}'^{\bullet},\theta'^{\bullet}|_H)\right)
$$
$$
=
\displaystyle
\bigoplus_{\lambda\in H^{\perp}}
\bigoplus_{p+q=\homdegree}
\prod_{i_0<\cdots<i_p}
\bigoplus_{\objectdegree\in\mathbb Z}
\Hom_{\Coh_G(U_{i_0\cdots i_p})}
\left(
(\mathcal F^{\objectdegree},\theta^{\objectdegree})
|_{U_{i_0\cdots i_p}}(\lambda),
(\mathcal F'^{\objectdegree+q},\theta'^{\objectdegree+q}|_H)
|_{U_{i_0\cdots i_p}}
\right).
$$
因此，作为 $\field$-向量空间的上链复形，有
$$
\hom_{\Coh_H(X)_{\dg}}
\left((\mathcal{F}^{\bullet},\theta^{\bullet}|_H),(\mathcal{F}'^{\bullet},\theta'^{\bullet}|_H)\right)
=
\hom_{\Coh_G(X)_{\dg}/H^{\perp}}
\left((\mathcal{F}^{\bullet},\theta^{\bullet}|_H),(\mathcal{F}'^{\bullet},\theta'^{\bullet}|_H)\right).
$$
\end{proof}

\begin{conceptbox}{轨道范畴（回顾）}
前面 \S3c 已给出轨道范畴的精确定义：\textbf{对象与 $\mathcal{C}$ 相同}，态射取
$\hom_{\mathcal{C}/K}(X,Y)=\bigoplus_{\lambda\in\widehat{K}}\hom_{\mathcal{C}}(X(\lambda),Y)$，
其中 $X(\lambda)$ 表示把等变结构按特征标 $\lambda$ 扭一下。本节用它把「先限制到子群 $H$、再商掉 $H^{\perp}$」写成全忠实嵌入，这正是扭 Künneth 公式在代数侧的来源；详细讨论与例子见 \S3c 的概念框。
\par\smallskip\noindent{\footnotesize 参考：Keller, \emph{On triangulated orbit categories}, arXiv:math/0610212。}
\end{conceptbox}
